% Background and Related Work: about 12 pages of the main body.
% Every citation is an entry of refs.bib verified on 2026-10-16 (refs_check.md). Statements about a cited work
% stay within what its abstract or official page says. Written in Stage 4.
\chapter{Background and Related Work}\label{ch02}

This chapter introduces the concepts the thesis builds on and the work it relates to. It follows the path of
the system: how knowledge graphs are modelled (\cref{sec:ch02-knowledge-graphs}), how knowledge is extracted
from tables and text (\cref{sec:ch02-knowledge-extraction-from-tables-a}), how the things they mention are
identified and linked (\cref{sec:ch02-entity-resolution-and-linking}), and how provenance is kept
(\cref{sec:ch02-provenance-and-traceability}); then how a graph and its texts are searched
(\cref{sec:ch02-information-retrieval,sec:ch02-retrieval-augmented-generation-and}), how language models act
through tools (\cref{sec:ch02-llm-agents-and-tool-use}) and query structured data
(\cref{sec:ch02-text-to-sql-and-question-answering-o}), and how such systems are evaluated
(\cref{sec:ch02-evaluation-with-llm-judges}). It closes by positioning this work against the closest systems
(\cref{sec:ch02-positioning-of-this-work}).

\section{Knowledge Graphs}\label{sec:ch02-knowledge-graphs}

A knowledge graph represents knowledge as a graph whose nodes stand for entities and whose edges stand for
relations between them~\cite{hogan2021knowledge}. Hogan et al.\ survey the area broadly: the data models and
query languages of graphs, schema, identity and context, deductive and inductive techniques, and the methods
used to create, enrich, assess, refine and publish knowledge graphs~\cite{hogan2021knowledge}. Ji et al.\
organise research on knowledge graphs into representation learning, knowledge acquisition and completion,
temporal graphs, and knowledge-aware applications~\cite{ji2022survey}. Large public graphs illustrate the
scale the field has reached: DBpedia extracts structured information from Wikipedia and publishes it as linked
data~\cite{auer2007dbpedia,lehmann2015dbpedia}, and Wikidata is a collaboratively edited knowledge base that
serves Wikipedia and other users~\cite{vrandecic2014wikidata}.

\paragraph{Data models.} Two families of graph data models dominate. In the RDF model, every fact is a
subject--predicate--object triple, and properties of a relation require additional constructs. In the
\emph{labelled property graph} model, used by graph databases such as Neo4j, nodes carry one or more labels,
edges carry a type, and both carry key--value properties. Graph database models have a long history of
data structures, query languages and integrity constraints~\cite{angles2008survey}. Schema support in property
graph systems remains limited; PG-Schema proposes a formalism for property graph schemas with type
definitions, inheritance and key constraints~\cite{angles2023pgschema}. This thesis uses the property graph
model, because provenance and identity decisions are naturally properties of edges: an edge from a mention
to the record it refers to can carry the reason, the score and the quote behind the decision.

\paragraph{Knowledge graphs and language models.} Pan et al.\ describe three ways of combining large language
models (LLMs) with knowledge graphs: knowledge graphs that enhance LLMs, LLMs that augment knowledge graphs
(including their construction), and systems in which both work together~\cite{pan2024unifying}. They argue
that graphs supply explicit and interpretable knowledge that LLMs lack, while graphs are hard to build and keep
current --- the motivation for using LLMs to build them. This thesis belongs to both of the latter groups: an
LLM helps to build the graph, and an LLM agent uses the graph to answer questions.

\section{Knowledge Extraction from Tables and Text}\label{sec:ch02-knowledge-extraction-from-tables-a}

\paragraph{From tables.} Turning tables into a graph requires deciding which columns identify an entity,
which rows become nodes and which become relations, and how the columns of different tables correspond. The
last of these is the schema-matching problem, surveyed by Rahm and Bernstein, who classify matchers by whether
they use the schema or the instances, single elements or structures, and names or constraints~\cite{rahm2001survey}.
In this thesis these decisions are proposed by an LLM as a declarative plan and validated by code, following
the pattern of the course that inspired the project, in which agents propose schemas for structured and
unstructured data and then build and connect both graphs~\cite{kollegger2025agentic}.

\paragraph{From text.} Open Information Extraction introduced the extraction of relational tuples for an
unbounded set of relations in one pass over a corpus, without hand-written rules or labelled examples for each
relation~\cite{banko2007open}. Generative LLMs have since become a general tool for information extraction;
Xu et al.\ survey this work by subtask (entity, relation and event extraction) and by technique~\cite{xu2024large}.
Zhu et al.\ evaluate LLMs on knowledge-graph construction and reasoning tasks and find GPT-4 better suited as
an inference assistant than as a few-shot information extractor~\cite{zhu2024llms}. Extract--Define--Canonicalize
(EDC) separates open extraction, the definition of a schema and the canonicalisation of the extracted
relations, and can build its own schema when none is given~\cite{zhang2024extract}. The present work shares
EDC's view that the schema need not be fixed in advance, but it differs in two respects: the schema proposed
for the text is validated by code and then fixed for the extraction, and every extracted item must carry a
quote that code finds verbatim in its source.

\paragraph{Frameworks.} Several open-source frameworks build graphs from documents with LLMs. The Neo4j LLM
Knowledge Graph Builder turns unstructured content into a graph of documents, chunks and extracted entities,
and supports vector, graph and text-to-Cypher querying; its documentation notes that it is less suited to
tabular data such as CSV files~\cite{neo4j2026llmgraphbuilder}. The \texttt{neo4j-graphrag} Python package
offers a knowledge-graph builder pipeline and retrievers over Neo4j vector indexes~\cite{neo4j2026graphrag}.
LlamaIndex's property graph index extracts a labelled property graph from text chunks with pluggable extractors,
from free triples to a strict schema, and offers retrievers that include text-to-Cypher~\cite{llamaindex2026propertygraph}.
These frameworks are built around text; the combination this thesis addresses --- records from several tables
and documents that talk about those records, built into one graph --- is outside their main use case.

\section{Entity Resolution and Linking}\label{sec:ch02-entity-resolution-and-linking}

\paragraph{Entity resolution} --- also called record linkage, data matching or duplicate detection --- identifies
records that refer to the same real-world entity, within one source or across several~\cite{christen2012data}.
Its probabilistic foundation is the model of Fellegi and Sunter, in which a vector of agreements between two
records leads to one of three decisions: link, possible link, or non-link, with the error rates of both kinds
held at chosen levels~\cite{fellegi1969theory}. The middle decision, a possible link sent for review, has a
direct counterpart in this work: pairs that rules cannot decide are sent to an LLM, whose decision code then
verifies. Since comparing every pair is quadratic, practical systems first restrict the candidates by
\emph{blocking}, comparing only pairs likely to match, or by \emph{filtering} on a similarity
threshold~\cite{papadakis2020blocking}; this work blocks by mutual nearest neighbours in an embedding space.

\paragraph{Entity linking} maps mentions in text to the entities of a knowledge base. Shen et al.\ identify its
two core difficulties as name variation (one entity under several names) and ambiguity (one name for several
entities)~\cite{shen2015entity}. Both appear in this thesis's data: one person written as ``Jonathan'', ``Jon''
and ``J.~Pike''; two people with the same name in different departments. Unlike classical entity linking, the
``knowledge base'' here is built in the same run from the dataset's own tables, and many mentions refer to
things no table holds, which must be kept as entities of their own rather than forced onto the nearest record.

\paragraph{Merging versus linking.} A common design merges duplicates into one node. An alternative keeps each
source item and records identity as a link between them. The second is reversible and keeps the evidence for
each decision; it is the design this thesis adopts and evaluates (\cref{ch04}).

\section{Provenance and Traceability}\label{sec:ch02-provenance-and-traceability}

Provenance records how a piece of data came to be. The W3C PROV data model describes it in terms of entities,
the activities that generate and use them, and the agents responsible, with relations such as
\emph{wasGeneratedBy}, \emph{used}, \emph{wasDerivedFrom} and \emph{wasAttributedTo}~\cite{moreau2013provdm}.
For a graph built by a language model, provenance answers a practical question: can a reader check, for any
node or edge, which source text or record supports it? Lewis et al.\ already named provenance --- the
difficulty of tracing what a purely parametric model knows to a source --- as a motivation for
retrieval-augmented generation~\cite{lewis2020retrieval}. This thesis applies the idea at two levels: every
item the builder derives from text carries a verbatim quote from its chunk, which code checks; and every answer
of the query engine cites the records and passages its tools returned, which code checks as well. The thesis
does not implement the PROV vocabulary itself; it records provenance as properties and edges of the property
graph and of the answer.

\section{Information Retrieval}\label{sec:ch02-information-retrieval}

\paragraph{Lexical retrieval.} BM25 ranks documents by a function derived from the Probabilistic Relevance
Framework, combining term frequency, inverse document frequency and document length~\cite{robertson2009probabilistic}.
It is fast, needs no training, and matches exact words and identifiers well, but it does not match
paraphrases. Lucene, the engine behind Neo4j's full-text indexes used in this work, implements BM25 ranking.

\paragraph{Dense retrieval} encodes questions and passages as vectors and ranks by similarity. Karpukhin et
al.\ showed that a dual encoder trained on questions and passages outperforms a strong BM25 system in top-20
retrieval accuracy for open-domain question answering~\cite{karpukhin2020dense}. M3-Embedding (BGE-M3) is an
example of a recent multilingual embedding model supporting dense, multi-vector and sparse retrieval in one
model~\cite{chen2024m3}. Dense retrieval matches meaning but can miss rare names and exact numbers --- the
mirror image of BM25's weakness.

\paragraph{Hybrid retrieval and fusion.} Because the two families fail differently, they are often combined.
Reciprocal Rank Fusion (RRF) scores each item by the sum, over the input rankings, of $1/(k + r)$, where $r$ is
the item's rank in that ranking and $k$ a constant~\cite{cormack2009reciprocal}:
\[
\mathrm{RRF}(d) = \sum_{i} \frac{1}{k + r_i(d)}.
\]
Cormack et al.\ found that RRF consistently outperformed the individual systems and Condorcet fusion on TREC
runs~\cite{cormack2009reciprocal}; it needs no score normalisation and no training. Bruch et al.\ give the main
counterpoint: in their analysis RRF is sensitive to its parameters, and a convex combination of normalised
scores performs better and can be tuned with few examples~\cite{bruch2023analysis}. This thesis uses RRF and
reports the setting of $k$ chosen for each use, together with the experiments that chose it (\cref{ch07}).

\paragraph{Reranking.} A first-stage retriever is often followed by a more expensive reranker applied to its
top candidates. A cross-encoder reads the question and a passage together; Nogueira and Cho's BERT re-ranker
set the state of the art on MS~MARCO passage ranking in this way~\cite{nogueira2019passage}. BGE-reranker-v2-m3
is a multilingual cross-encoder of this kind~\cite{baai2024bgererankerv2m3}. Generative LLMs can also rerank:
Sun et al.\ showed that suitably instructed LLMs used as zero-shot rerankers match or exceed supervised methods
on standard benchmarks~\cite{sun2023chatgpt}. \Cref{ch07} compares an LLM reranker with a cross-encoder on this
thesis's data.

\section{Retrieval-Augmented Generation and GraphRAG}\label{sec:ch02-retrieval-augmented-generation-and}

\paragraph{RAG.} Retrieval-augmented generation pairs a generator with a retriever over an external
index~\cite{lewis2020retrieval}: the model answers from passages retrieved for the question rather than from its
parameters alone. Gao et al.\ distinguish naive, advanced and modular forms of RAG~\cite{gao2023retrieval}.
Two known limits of plain RAG matter here. First, questions that collect or count every match across a corpus
cannot be answered from a few retrieved passages. Second, even relevant passages may be under-used: Liu et al.\
show that LLMs use information at the beginning or end of a long context better than information in the
middle~\cite{liu2024lost}.

\paragraph{GraphRAG.} Graph-based RAG uses a graph to index, retrieve or organise the material given to the
generator. Peng et al.\ describe its workflow in three stages: graph-based indexing, graph-guided retrieval and
graph-enhanced generation~\cite{peng2026graph}. Two systems are closest to this work. Microsoft's GraphRAG has
an LLM build an entity graph from the documents and summarise communities of related entities in advance; it
targets ``global'' questions about a whole corpus, which it answers by combining partial answers from the
community summaries~\cite{edge2024local}. HippoRAG has an LLM extract a graph from passages and retrieves
passages for multi-hop questions with Personalized PageRank over that graph~\cite{gutierrez2024hipporag}.
In both systems the graph is derived from text only. A further, practical pattern combines a \emph{lexical}
graph of documents and chunks with a \emph{domain} graph of entities, as in Neo4j's GraphRAG
tooling~\cite{neo4j2026graphrag,neo4j2026llmgraphbuilder}.

The value of a graph over structured data has also been measured directly: on an enterprise insurance
database, Sequeda et al.\ report that zero-shot GPT-4 answered 16\,\% of questions correctly when querying
the SQL schema directly but 54\,\% when querying a knowledge-graph representation of the same
data~\cite{sequeda2024benchmark}. This thesis studies a different question: not whether a graph helps an LLM
to write queries, but what the graph should contain when tables and text must be answered together. Its
central finding (\cref{ch04,ch05}) is that, for this purpose, the graph is more useful as an index of
identity and evidence than as a store of every extracted fact.

\section{LLM Agents and Tool Use}\label{sec:ch02-llm-agents-and-tool-use}

An LLM agent is a model that decides, step by step, which actions to take towards a goal. Wang et al.\ survey
such agents and propose a common framework for their construction~\cite{wang2024survey}. Chain-of-thought
prompting showed that asking a model to write intermediate reasoning steps improves multi-step
reasoning~\cite{wei2022chain}. ReAct interleaves such reasoning with actions that query external sources, and
reduced hallucination and error propagation compared with reasoning alone~\cite{yao2023react}. Toolformer
trains a model to decide which tool to call, when, with which arguments, and how to use the
result~\cite{schick2023toolformer}. Modern LLM APIs expose the same capability as \emph{function calling}, and
frameworks such as Google's Agent Development Kit package it for building agents~\cite{google2026adk}.

\paragraph{Agents over graphs and structured data.} Think-on-Graph treats the LLM as an agent that explores a
knowledge graph by iterative beam search over entities and relations~\cite{sun2024thinkongraph}. Reasoning on
Graphs first plans relation paths grounded in the graph, then retrieves and reasons over the matching
paths~\cite{luo2024reasoning}. StructGPT gives the model specialised interfaces for reading knowledge graphs,
tables and databases, and alternates reading with reasoning~\cite{jiang2023structgpt}. The query engine of this
thesis belongs to this family: its agent calls typed, read-only tools over the graph, the passages and the
tables. It differs in that its graph was built automatically from the same mixed data, and in that code checks
every tool call, every citation and the form of every answer before an answer is returned.

\section{Text-to-SQL and Question Answering over Tables and Text}\label{sec:ch02-text-to-sql-and-question-answering-o}

\paragraph{Text-to-SQL and text-to-Cypher.} Translating questions into database queries is a long-standing
task. Spider is a cross-domain text-to-SQL benchmark with disjoint training and test databases, which requires
generalising to unseen schemas~\cite{yu2018spider}. BIRD stresses database values, dirty contents and external
knowledge, and reported a large gap between ChatGPT and humans in execution accuracy~\cite{li2023bird}. For graph
databases, the Text2Cypher dataset collects natural-language questions with Cypher queries and shows that
fine-tuned models improve on baselines~\cite{ozsoy2025text2cypher}. Free query generation is powerful but must be
constrained in a system that answers arbitrary users: a generated query may be wrong, expensive, or attempt to
change the data. This thesis's early query stage used text-to-Cypher with checks; its final engine instead
gives the agent a typed query language that code compiles to Cypher or SQL, and evaluates free SQL under
validation as an alternative (\cref{ch06,ch07}).

\paragraph{Tables and text together.} Benchmarks for questions that need both tables and text show the
difficulty of the combination. In HybridQA, each question needs a Wikipedia table and linked passages; models
using only the table or only the text scored far below a hybrid model, which in turn scored far below
humans~\cite{chen2020hybridqa}. TAT-QA, built from financial reports, adds numerical reasoning over tables and
text~\cite{zhu2021tatqa}. These benchmarks provide the table and its linked passages with each question; the
setting of this thesis is harder in one respect --- the links between text and records are not given, and the
builder must find them --- and easier in another, since its corpora are small.

\section{Evaluation with LLM Judges}\label{sec:ch02-evaluation-with-llm-judges}

Evaluating free-text answers by string match undercounts correct paraphrases, while human judging is
expensive. Strong LLMs are therefore used as judges. Zheng et al.\ found that GPT-4's judgements agreed with
human preferences more than 80\,\% of the time, about as often as humans agree with each other, but also
identified position, verbosity and self-enhancement biases~\cite{zheng2023judging}. Self-enhancement --- a
model's tendency to favour its own outputs --- is why this thesis never lets the model that built a graph or
answered a question judge it. RAGAS evaluates RAG pipelines without reference answers, along dimensions such as
the relevance of the retrieved context and the faithfulness of the answer to it~\cite{es2024ragas}. This thesis
uses reference answers instead, and reports a deterministic exact-match score next to every judged score
(\cref{sec:ch03-scoring-exact-match-and-judgement}).

For comparing two systems on the same questions, Dietterich's analysis of statistical tests for comparing
learning algorithms recommends McNemar's test, which has a low type~I error, over tests that treat the two
systems' results as independent~\cite{dietterich1998approximate,mcnemar1947note}. Rates are reported with
Wilson score intervals~\cite{wilson1927probable}; both are defined in \cref{sec:ch03-statistical-analysis}.

\section{Positioning of This Work}\label{sec:ch02-positioning-of-this-work}

\Cref{tab:positioning} compares this work with the closest systems along the dimensions that define it. The
comparison describes what each system is designed for, as stated by its authors; it is not a measured
comparison, since the systems target different tasks and data.

\begin{table}[htbp]
\centering
\caption{This work and the closest related systems, by design.}
\label{tab:positioning}
\scriptsize
\renewcommand{\arraystretch}{1.25}
\begin{tabularx}{\linewidth}{@{}>{\raggedright}p{2.1cm}>{\raggedright}X>{\raggedright}X>{\raggedright\arraybackslash}X@{}}
\toprule
System & Input and graph & Role of the graph & Querying \\
\midrule
GraphRAG~\cite{edge2024local} & text; LLM-extracted entity graph with community summaries & summarise a whole corpus & combine partial answers from community summaries \\
HippoRAG~\cite{gutierrez2024hipporag} & text; LLM-extracted graph over passages & find passages for multi-hop questions & Personalized PageRank, then read \\
LLM Graph Builder~\cite{neo4j2026llmgraphbuilder} & mainly unstructured content; chunk graph plus entity graph & store and navigate extracted entities & vector, graph and text-to-Cypher \\
LlamaIndex property graph~\cite{llamaindex2026propertygraph} & text chunks; extracted property graph & store extracted paths & synonym, vector, text-to-Cypher retrievers \\
EDC~\cite{zhang2024extract} & text; extracted, then canonicalised triples & the output itself & --- \\
Think-on-Graph~\cite{sun2024thinkongraph} & an existing knowledge graph & reasoning paths & LLM beam search over the graph \\
StructGPT~\cite{jiang2023structgpt} & existing graphs, tables, databases & one of several sources & read through interfaces, then reason \\
\midrule
\textbf{This work} & \textbf{tables and text together; records from an LLM plan checked by code; mentions linked to records by reversible, evidenced edges} & \textbf{an index of identity and evidence, with a stated navigation contract} & \textbf{agent with typed, read-only graph, text and table tools; code-checked citations and outcomes} \\
\bottomrule
\end{tabularx}
\end{table}

Three points distinguish this work. First, it builds one graph from \emph{both} tables and text in a single
automated pipeline, and treats the links between them --- which text mentions which record --- as the part that
must be exact. Second, it measures what the graph should contain: it builds and evaluates a sequence of
representations, from a graph of extracted facts to an \emph{anchor graph} that indexes identity and evidence,
and evaluates the graph as an index with criteria fixed in advance rather than only through the answers of a
question-answering system. Third, its query engine separates what a language model decides from what code
verifies, at every step from the query to the citation.

The techniques it uses are established: BM25, dense retrieval, RRF, LLM and cross-encoder reranking, the ReAct
pattern of interleaved reasoning and tool calls, blocking for entity resolution, and the statistical tests
above. The contribution lies in how they are combined and constrained for mixed tables and text, in the
evidence gathered on which representation serves this purpose, and in the evaluation method itself
(\cref{sec:ch01-contributions}).
